\subsection{LIMITS IN PRODUCT SPACES AND QUOTIENT SPACES}

PROPOSITION 10. Let X be a set, let \((Y_{i})_{i\in I}\) be a family of topological spaces, and for each \(i\in I\) let \(f_{i}\) be a mapping of X into \(Y_{i}\) . Let X be given the coarsest topology G for which the \(f_{i}\) are continuous. Then a necessary and sufficient condition for a filter F on X to converge to \(a\in X\) is that for each \(i\in I\) the filter base \(f_{i}(\mathfrak{F})\) should converge to \(f_{i}(a)\) in \(Y_{i}\) .

The condition is necessary since the \(f_{i}\) are continuous (no. 4, Proposition 9, Corollary 1). Conversely, suppose that the condition is satisfied, and let V be an open neighbourhood of \(a\) in X. By the definition of \(\mathfrak{C}\) ( \(\S\) 2, no. 3, Proposition 4) there is a finite subset J of I, and for each \(i\in J\) an open subset \(\mathrm{U}_{i}\) of \(\mathrm{Y}_{i}\) , such that \(f_{i}(a)\in\mathrm{U}_{i}\) for \(i\in J\) and such that V contains the set

\[
\bigcap_ {i \in J} \overline {{f}} _ {i} ^ {1} (U _ {i}).
\]

The hypothesis implies that \(\overline{f}_{i}^{1}(U_{i})\in\mathfrak{F}\) (no. 3, Proposition 7); since J is finite, it follows that

\[
\mathbf {M} = \bigcap_ {i \in \mathbf {J}} \overline {{f}} _ {i} ^ {1} (\mathrm{U} _ {i})
\]

belongs to \(\mathfrak{F}\) , and \(\mathbf{M} \subset \mathbf{V}\) . This completes the proof.

COROLLARY 1. A filter \(\mathfrak{F}\) on a product space \(\mathbf{X} = \prod_{i\in I}\mathbf{X}_i\) converges to a point \(x\) if and only if for each \(i\in I\) the filter base \(\operatorname{pr}_{i}(\mathfrak{F})\) converges to \(\operatorname{pr}_{i}(x)\) .

COROLLARY 2. Let \(f = (f_{i})\) be a mapping of a set X into a product space \(Y = \prod_{i \in I} Y_{i}\) . Then \(f\) has a limit \(y = (y_{i})\) with respect to a filter \(\mathfrak{F}\) on X if and only if for each \(i \in I\) \(f_{i}\) has limit \(y_{i}\) with respect to \(\mathfrak{F}\) .

PROPOSITION 11. Let R be an open equivalence relation on a topological space X and let \(\varphi\) be the canonical mapping \(X \rightarrow X/R\) . Then for each \(x \in X\) and each filter base \(B'\) on X/R which converges to \(\varphi(x)\) , there is a filter base B on X which converges to x and is such that \(\varphi(\mathfrak{B})\) is equivalent to \(B'\) .

If U is any neighbourhood of x in X, then \(\varphi(\mathbf{U})\) is a neighbourhood of \(\varphi(x)\) in X/R (§ 5, no. 3, Proposition 5), hence there is a set \(M' \in B'\) such that \(M' \subset \varphi(\mathbf{U})\) ; if we put \(M = U \cap \overline{\varphi}^{1}(M')\) , then \(M' = \varphi(M)\) . This shows that as \(M'\) runs through \(B'\) and U runs through the neighbourhood filter of x, the sets \(U \cap \overline{\varphi}^{1}(M')\) form a filter base B on X; clearly B converges to x and \(\varphi(\mathfrak{B})\) is equivalent to \(B'\) .

